\section{Background and Problem Formulation}
\label{sec:background}

\subsection{Affected-Version Identification}

This paper studies the vulnerability-affected version identification task defined by the benchmark study of Chen et al.~\cite{howfar2025}. Given a reported vulnerability $c$, such as a CVE, a target project or repository $P$, a set of released versions or release tags $V=\{v_1,\ldots,v_n\}$, and patch evidence such as one or more fixing commits, the task is to identify the subset of releases in which the vulnerable logic remains present:
\[
V_{\mathit{aff}}(c,P) = \{v \in V \mid \mathit{vuln\_state}(c,P,v)=\mathsf{true}\}.
\]
Here $\mathit{vuln\_state}$ is a task-level predicate, not an assumed oracle. It names the conceptual release-level property that tag $v$ still contains the vulnerability condition associated with $c$. Later sections operationalize this property through repository evidence and branch-specific history reasoning.

The output is therefore a set of released versions, not a vulnerability-inducing commit, not a vulnerable function, and not a textual advisory range. This distinction is central to the benchmark study by Chen et al.: the evaluation treats exact recovery of the full affected-version set as a vulnerability-level objective, while also measuring partial correctness over individual CVE-version pairs~\cite{howfar2025}. A method can be useful at one granularity and still fail at the other. For example, predicting many truly affected tags may yield a reasonable version-level recall, but a single omitted affected tag makes the CVE-level exact set wrong.

The benchmark setting is also repository-grounded. The released-version universe is determined by project release tags, and the ground-truth labels are attached to those released versions. In this paper, we use the same task framing and benchmark family as Chen et al.: the goal is to predict \texttt{affected\_versions} for each CVE in the dataset, using repository and patch evidence rather than treating public advisory metadata as sufficient ground truth.

\subsection{From Fixing Patches to Vulnerability State}

A fixing patch is the most useful starting signal for this task, but it is not the answer. A commit that closes a CVE may contain the root-cause fix, supporting checks, error-handling changes, tests, documentation, refactoring, wrapper changes, formatting edits, or branch-specific backport adjustments. The patch also describes a transition from one repository state to another, whereas affected-version identification asks whether many released states still satisfy the vulnerable condition.

We use two conceptual predicates to separate these roles. A \emph{root-cause predicate} describes the condition that makes the vulnerability possible in a repository state, such as a missing bounds check, an unchecked object lifetime, an attacker-controlled value reaching a sensitive operation, or an invalid state transition. A \emph{fix predicate} describes the condition introduced or restored by the patch that prevents the vulnerability. Code presence alone may also be insufficient when the vulnerability depends on attacker-controlled inputs, trigger conditions, configuration, reachable sinks, or exploit preconditions. These attacker-relevant conditions should be preserved as structured evidence rather than only as unstructured prompt text. These predicates are conceptual objects in this section; the concrete extraction and checking procedures are deferred to Section~\ref{sec:approach}.

This predicate view prevents a common category error. The old side of a patch may provide a pre-fix anchor, but the anchor is only evidence about where to search. The new side of a patch may provide a fix predicate, but the absence of that new code in an old release does not automatically prove that the old release is vulnerable. A release is affected only when the root-cause predicate, under the relevant repository context, still holds in that release.

\subsection{Attacker Conditions and Vulnerability Evidence Graphs}

Affected-version identification cannot be reduced to code similarity, patch-line matching, or BIC localization. A vulnerable line matters only when the surrounding repository state still permits the vulnerable behavior. This state often depends on attacker-relevant conditions: attacker-controlled input, a trigger condition, an exploit precondition, a reachable sink, a missing guard, an unsafe state transition, or a security condition such as object lifetime, bounds, recursion, integer overflow, or a type invariant.

For this reason, a root-cause predicate is not just a code line. It should be tied to the vulnerable condition, the fix predicate or fix condition, source evidence, patch evidence, and the path or state constraint that makes the vulnerability relevant. The evidence may include a reachable sink, a missing guard, an unsafe state transition, a pre-fix anchor, or a history event that preserves or changes the condition. For affected-version identification, an attacker perspective does not mean proving a working exploit for every release. It means preserving the conditions under which the vulnerable behavior can be triggered, so that later history reconstruction does not confuse a syntactically similar line with a security-relevant state.

These relations are hard to preserve as flat prompt text. They are also easy to lose in a flat JSON record that lists evidence items without their roles. The evidence for one CVE can connect a CVE description, CWE category, fixing commit, patch hunk, pre-fix anchor, root-cause predicate, fix predicate, attacker condition, history event, and boundary event. A graph representation provides a structured substrate for these relations. In this limited sense, an attack-perspective knowledge graph is a way to organize attacker conditions and repository evidence, not a generic security ontology or an exploit proof. Its purpose is narrower: to keep evidence connected across the affected-version pipeline, from root-cause evidence to SZZ anchors, history events, boundary-event judgment, Git-DAG state propagation, and release tag projection.

\subsection{Why BIC Localization Is Insufficient}

SZZ-style methods locate bug-inducing commits (BICs) or vulnerability-inducing commits (VICs) by tracing from a fixing commit back through file history. This evidence is useful because the introduction of a vulnerable statement or state can bound where the vulnerability may have appeared. V-SZZ connects this idea to affected-version inference by mapping vulnerability-inducing commits to released version ranges~\cite{vszz2022}; later work improves commit localization through semantic analysis, LLM-assisted statement and candidate assessment, agentic repository search, and temporal knowledge-graph search~\cite{tcszz2024,llm4szz2025,agentszz2026,howwhyagents2026,beyondblame2026}.

However, a BIC or VIC is a history event, while an affected-version set is a projection over released repository states. Even a correct inducing commit does not by itself determine all affected releases. The method must still determine which branch-specific states reach official release tags, which branches received the fix, whether equivalent fixes exist without the same commit hash, and whether later code movement or partial reverts changed the vulnerability state.

The gap becomes larger in non-linear histories. A vulnerability may be introduced on one branch, fixed on another, backported to selected stable branches, or fixed through different commits in different release lines. Multi-fix CVEs and partial fixes further break the assumption that a single inducing commit and a single fixing commit define a closed linear interval. Therefore, BIC localization provides history evidence, but affected-version identification requires boundary-event judgment and branch-aware release projection.

\subsection{Tracing-Based and Matching-Based Evidence}

The benchmark study organizes existing affected-version tools into two evidence families~\cite{howfar2025}. Tracing-based methods start from fixing patches, select statements or hunks, trace them through history toward candidate inducing commits, and then infer affected versions from the resulting history range. Matching-based methods construct vulnerability signatures, patch signatures, or code representations and search historical versions for matching vulnerable logic.

Both families produce useful evidence, but neither evidence family is equivalent to complete affected-version reasoning. Tracing evidence can identify plausible historical events, yet it may be sensitive to statement selection, one-step blame assumptions, over-tracing, and weak cross-branch patch reuse. Matching evidence can detect code presence, yet it may be brittle under renaming, refactoring, minor syntactic edits, and semantically equivalent but syntactically different code.

The benchmark evaluation deliberately uses both vulnerability-level and version-level metrics because the two evidence families fail in different ways~\cite{howfar2025}. Tracing-based tools often preserve temporal context but can miss the true vulnerable condition when the selected patch lines are not root-cause lines. Matching-based tools can achieve high precision on exact or near-exact patterns but often lack semantic verification that a matched snippet still constitutes the vulnerability. For this reason, the present paper treats tracing and matching outputs as evidence sources rather than final affected-version decisions.

\subsection{History and Release Structure Challenges}

Affected-version identification is difficult because patch structure and release structure interact. Add-only patches are hard for deletion-dependent tracing because the patch introduces guards or checks without deleting the vulnerable old statement. Deleted-line or modified-line patches are easier to anchor, but the deleted text may still be an implementation symptom rather than the root cause. Multi-function and multi-file patches add unrelated edits, helper changes, tests, or refactoring noise, making it harder to decide which code elements should drive the affected-version judgment.

Branch and release structure adds another layer. Many projects maintain multiple release lines, and a CVE fix may be applied to each line through cherry-picks, backports, merge commits, or semantically equivalent commits. A branch may receive only part of a fix, may revert a fix, or may already contain a different mitigation. Consequently, the same CVE can correspond to several boundary events: the first commit that introduces the vulnerable state, the fixing commit on one branch, equivalent fixes on other branches, and prerequisite commits that make the vulnerability reachable.

Code evolution also weakens naive signature reasoning. A vulnerable function can be renamed, split, copied, moved across files, wrapped by a new API, or refactored while preserving the same vulnerable behavior. Conversely, a syntactically similar function can be safe because a guard, configuration constraint, type invariant, or control-flow condition changed. The release universe itself must also be explicit. Affected-version predictions are evaluated over official released versions or release tags, not arbitrary Git refs or every historical tag.

\subsection{Evidence-Gated Continual Adaptation}

Affected-version identification has repeated structure across CVEs. Cases from the same repository may share release branches, backport style, subsystem layout, parser code, or dependency code. Cases with the same CWE may share a root-cause pattern, such as bounds checks, object lifetime, recursion limits, integer overflow, or type confusion. Fixing commits may also share patch patterns, such as add-only guards, changed validation logic, changed error paths, or branch-specific backports.

These repeated structures can help later CVEs, but only as evidence. They can inform pre-fix anchor selection, candidate history event ranking, boundary-event judgment, branch-specific fix reasoning, and failure diagnosis. For example, a verified observation that one repository backports fixes through equivalent stable-branch commits can help rank future candidate fixes in that repository. A verified failure mode for add-only guard patches can help avoid relying only on deleted lines. In both cases, the reusable information is the inspected evidence and its provenance, not a model's unsupported memory of a prior answer.

Direct model memory is risky for this task. It may hallucinate a prior pattern, overfit to development cases, preserve a wrong prior, leak benchmark labels, or hide the source of an error. Continual adaptation is useful only if it is evidence-gated. A previous decision should not become reusable knowledge merely because a model produced it. It should become reusable only when the decision is linked to inspected repository evidence, a stable patch pattern, or a verified failure mode.

Evidence-gated memory should therefore keep provenance, link observations to repository, CWE, and patch-pattern evidence, and allow an observation to be invalidated or downgraded when later evidence contradicts it. It should also separate reusable evidence from ground-truth affected-version labels. In this paper, continual adaptation is a design requirement and a future ablation target unless supported by explicit experiments. It should act as a prior for history reconstruction, not as the final affected-version decision.

\subsection{Problem Statement}

The problem studied in this paper is: given a CVE $c$, its fixing commit family $F_c$, a target repository $P$, and the official release-tag universe $V$, identify all and only release tags in $V$ whose repository states satisfy the root-cause vulnerability predicate for $c$. Here, $F_c$ denotes one or more fixing commits associated with the CVE, including branch-specific or equivalent fixes when such evidence is available. The expected output is a set of affected release tags, together with evidence and uncertainty records that explain how the prediction was derived.

This formulation deliberately excludes several adjacent tasks. It is not exploit generation and does not attempt to prove runtime exploitability for every released tag. It is not merely advisory normalization because advisory ranges can be incomplete, stale, or inconsistent with repository history. It is not only BIC or VIC localization because inducing commits are history events rather than release-level exposure labels. It is also not only vulnerable-code clone detection because code similarity is evidence, not a vulnerability-state verdict.

We describe the repository history as a Git directed acyclic graph (Git DAG). A \emph{history event} is a commit-level change that may introduce, remove, preserve, or transform part of the vulnerability condition. A \emph{boundary event} is a history event that changes the vulnerability state for a branch or release line. A \emph{release tag projection} maps these branch-specific vulnerability states onto the official release-tag universe. These definitions provide the vocabulary needed by the later method section without committing this section to a particular algorithm.

\subsection{Requirements for a Robust Solution}

The preceding formulation implies that an effective solution must be root-cause-grounded. It must distinguish vulnerability-relevant evidence from incidental patch edits, and it must represent both the vulnerable condition and the fixing condition before projecting branch-specific states to release tags. Without this separation, a method can confuse a wrapper change, a test update, or a backport artifact with the actual security state.

The solution must also be evidence-constrained. Pre-fix anchors, history events, and boundary events should be accepted only with repository evidence that can be inspected later. In particular, the method should separate the judgment that a candidate history event is relevant from the projection that converts branch-specific states into affected release tags. This separation is necessary because a candidate commit can be useful evidence without being the final affected-version boundary.

Finally, the solution must be branch-specific and auditable. It should reason over the Git DAG, release lines, backports, equivalent fixes, and official release tags rather than assuming a single linear version order. If LLMs or agents are used for semantic judgment, their role should be bounded and accountable: prompts, evidence inputs, uncertainty, failures, and decision artifacts must be recorded so that affected-version predictions can be inspected, debugged, and evaluated under the benchmark metrics. A robust solution should also support evidence-gated accumulation of repeated failures and verified history-event decisions as reusable repository, CWE, and patch-pattern evidence. Such memory must be auditable, not free-form model memory.
